% Source: https://artofproblemsolving.com/community/c3h1357521_q_on_mass_points
% Stable identity: aops-topic-1357521-q01
\textbf{Mass points with a split mass.} Line $EF$ does not pass through a vertex, so we split the mass at $A$ into a part $a_1$ that balances $B$ along $AB$ and a part $a_2$ that balances $C$ along $AC$. Take $m_B=m_C=1$ (so $D$ is the balance point of $BC$). Then $AF:FB=1:2$ needs $a_1:m_B=2:1$, so $a_1=2$, and $AE:EC=1:3$ needs $a_2:m_C=3:1$, so $a_2=3$. The total mass at $A$ is $a_1+a_2=5$, and the mass at $D$ is $m_B+m_C=2$.

\begin{asy}
usepackage("amsmath");
size(5.5cm);
defaultpen(fontsize(10pt));
pair B=(0,0),C=(8,0),A=(3,6),D=(4,0),E=(4.25,4.5),F=(2,4),P=A+2/7*(D-A);
draw(A--B--C--cycle,linewidth(0.9));
draw(A--D,blue); draw(E--F,red);
dot(D);dot(E);dot(F);dot(P);
label("$A\ (5)$",A,N); label("$B\ (1)$",B,SW); label("$C\ (1)$",C,SE); label("$D\ (2)$",D,S); label("$E$",E,NE); label("$F$",F,W); label("$P$",P,0.8*SE);
label("\mbox{}",(min(currentpicture,user=true)+max(currentpicture,user=true))/2,black); // keeps the text after the figure black when LaTeX typesets the labels inline
\end{asy}

The point $P$ is the balance point of the whole system: it lies on $FE$ (masses $a_1+m_B=3$ at $F$ and $a_2+m_C=4$ at $E$) and on $AD$ (mass $5$ at $A$, $2$ at $D$). Hence
\[
AP:PD=m_D:m_A=\boxed{2:5}.
\]
The three ratios in the question are not in conflict: $m_A:m_B=2:1$ and $m_A:m_C=3:1$ refer to the two different parts of $A$'s mass.

\textbf{Check with vectors.} With $A$ as origin, $D=\dfrac{B+C}{2}$, $E=\dfrac{C}{4}$, $F=\dfrac{B}{3}$. A point $tD=\dfrac {t}{2}B+\dfrac {t}{2}C$ of $AD$ lies on $FE$ when it equals $(1-s)\dfrac {B}{3}+s\dfrac {C}{4}$, i.e.\ $\dfrac {t}{2}=\dfrac{1-s}{3}$ and $\dfrac {t}{2}=\dfrac {s}{4}$. Then $s=2t$ and $3t=2-4t$, so $t=\dfrac{2}{7}$: $AP:PD=2:5$.
